\chapter{Programmation linéaire à deux variables}
\label{workbook:chapter:13}
\section{Variables, contraintes et fonction objectif}
\label{workbook:section:13.01}
\begin{exerciseblock}{Variables, contraintes et fonction objectif}
\item Une entreprise produit des tables $x$ et des chaises $y$. Traduire : au plus $80$ heures, $4$ h/table et $2$ h/chaise; au moins $10$ tables; quantités non négatives.
\item Si le profit vaut $50\,$\$ par table et $30\,$\$ par chaise, écrire la fonction objectif.
\item Expliquer pourquoi les unités aident à vérifier chaque coefficient.
\end{exerciseblock}
\section{Région réalisable}
\label{workbook:section:13.02}
\begin{exerciseblock}{Région réalisable}
\item Tracer la région définie par $x\ge0$, $y\ge0$, $x+y\le6$, $2x+y\le8$.
\item Trouver tous les sommets de cette région.
\item Tester si $(3,3)$, $(4,0)$ et $(1,5)$ sont réalisables.
\end{exerciseblock}
\section{Principe d'optimalité aux sommets}
\label{workbook:section:13.03}
\begin{exerciseblock}{Principe d'optimalité aux sommets}
\item Évaluer $P=3x+2y$ aux sommets de la région précédente.
\item Expliquer géométriquement pourquoi une fonction linéaire atteint un optimum à un sommet lorsqu'un optimum fini existe sur un polygone.
\item Donner un cas d'optimum atteint sur tout un segment.
\end{exerciseblock}
\section{Configurations particulières du domaine réalisable}
\label{workbook:section:13.04}
\begin{exerciseblock}{Configurations particulières du domaine réalisable}
\item Donner un exemple de région vide avec deux inégalités contradictoires.
\item Donner une région non bornée pour laquelle un maximum de $x+y$ n'existe pas.
\item Construire une contrainte redondante et expliquer comment la reconnaître.
\end{exerciseblock}
\section{Optimiser en faisant glisser une droite}
\label{workbook:section:13.05}
\begin{exerciseblock}{Optimiser en faisant glisser une droite}
\item Pour $P=4x+y$, dessiner trois lignes de niveau $P=4,8,12$ et indiquer leur direction de déplacement pour augmenter $P$.
\item Sur la région $x,y\ge0$, $x+2y\le8$, $3x+y\le9$, trouver le maximum graphiquement puis par les sommets.
\item Relier le vecteur $(4,1)$ à la direction de croissance de l'objectif.
\end{exerciseblock}
\section{De la phrase à la région réalisable}
\label{workbook:section:13.06}
\begin{exerciseblock}{De la phrase à la région réalisable}
\item Traduire : « au moins deux fois plus de $x$ que de $y$ », « total inférieur à $50$ », « $y$ entre $5$ et $20$ ».
\item Pour chaque inégalité, choisir un point test afin de déterminer le bon demi-plan.
\item Expliquer pourquoi $x,y\ge0$ ne doit pas être oublié dans un problème de production.
\end{exerciseblock}
\section{De la décision au modèle géométrique}
\label{workbook:section:13.07}
\begin{exerciseblock}{De la décision au modèle géométrique}
\item Formuler entièrement un problème de mélange de deux produits avec coût minimal et deux contraintes nutritionnelles.
\item À partir d'un tableau de coefficients, nommer variables, unités, contraintes et objectif.
\item Décrire une procédure de validation de la solution optimale dans le contexte initial.
\end{exerciseblock}
\section*{Synthèse du chapitre}
\begin{synthesisbox}
\item Un atelier fabrique $x$ produits A et $y$ produits B. Contraintes : $2x+y\le100$, $x+3y\le120$, $x,y\ge0$. Profit $40x+50y$. Résoudre complètement.
\item Créer un problème où l'optimum est multiple sur une arête et le démontrer par les lignes de niveau.
\item Analyser ce qui change si les variables doivent être entières alors que le sommet optimal a des coordonnées non entières.
\end{synthesisbox}
